\datos{color}{}{}
\section*{Instalación y despliegue}
\addcontentsline{toc}{section}{Instalación y despliegue}

Con el objetivo de facilitar la instalación y el despliegue del entorno 
de desarrollo utilizado, se recomienda utilizar la imagen \textit{Docker}
que se ha preparado. Esta imagen incluye tanto el simulador 
\textit{CoppeliaSim} como el ecocsistema de \textit{ROS2}. \\

Se puede construir utilizando el fichero \textit{Dockerfile}, 
Código \ref{cod:dockerfile}, con el comando ``\texttt{docker build -t coppelia\_sim\_ros2:jazzy .}''
o descargar mediante el comando: ``\texttt{docker pull arambarricalvoj/coppelia\_sim\_ros2:jazzy}''. \\

Una vez descargada la imagen, para crear un contenedor o instancia 
de la misma ejecutar desde terminal el fichero \textit{run.sh},
Código \ref{cod:run.sh}. La terminal accederá al contenedor \textit{Docker}
iniciado y arrancará automáticamente el simulador. \\

Para ejecutar la aplicación programada en \textit{ROS2}, desde 
otra terminal ejecutar ``\texttt{docker exec -it coppelia\_sim bash}''. 
La aplicación estará en la ruta 
\texttt{/home/\$USER/pcr\_misc/tareas/ros2\_ws}. \\

Todo este proceso junto con los comandos específicos para ejecutar
cada tarea se detalla en más profundidad en el repositorio del 
proyecto: \url{https://github.com/arambarricalvoj/pcr_misc_ucm-uned_25-26}.